\SourceLeaf{284}%
\begin{center}\begin{tabular}{@{}r c@{}}
8.1 & $-$\\
9.2 & $+$\\
12.5 & $-$\\
14.7 & $+$\\
16.11 & $-$\\
17.4 & $-$\\
18.3 & $+$\\
19.10 & $+$\\
20.15 & $+$\\
21.6 & $-$\\
22.13 & $-$\\
\end{tabular}\end{center}

Es ist klar, dass sich, wenn man von einem Punkt vor (1) ausgeht und auf der
Kurve die Punkte 1, 2, 3 etc. durchläuft, zum ersten Mal beim
Schnitt 8 eine Figur bildet, die wir $a$ nennen wollen. Sie geht von 1
bis 8. Offenbar liegen der Teil der ganzen Linie unmittelbar vor 1 und
[der Teil] nach 8 auf der $-$-Seite von $a$.

Geht man so weiter, bildet sich eine neue zusammenhängende
Kurve, die vom Anfangspunkt bis 1 und (unter Verlassen des vorigen Wegs 2, 3 etc.)
dann weiter über 9, 10, 11 etc. verläuft. Bei 9 gelangt man auf die
$+$-Seite von $a$, bei 12 auf die $-$-Seite, bei 14 auf die $+$-Seite von $a$.
Bei 16 hingegen trifft die Kurve sich selbst und schließt eine zweite
Figur, die von 11 über 12, 13, 14, 15 bis 16 geht. Sie heiße $b$.
Offenbar liegen der Teil der Kurve unmittelbar vor 1, dann das ganze
Stück von 8 bis 11 und schließlich der unmittelbar an 16 grenzende Teil auf
der $-$-Seite von $b$. Nachdem jetzt auch $b$ abgesondert ist, trifft die
Kurve bei 17 wieder $a$ und kommt auf die $-$-Seite, bei 18 trifft sie erneut darauf und
gelangt auf die $+$-Seite. Bei 19 schließt sich die $3^{\mathrm{te}}$
Figur (da 10 noch nicht in einem ausgeschlossenen Teil liegt), die
$c$ heißen soll und von 10 bis 11 $=$ 16, dann von 16 bis 19 geht.
Wird diese wiederum abgesondert, so liegt das verbleibende Stück
ebenso wie das unmittelbar folgende auf der $+$-Seite von $c$. Bei 20 gelangt
man auf die $+$-Seite von $b$, bei 21 auf die $-$-Seite von $a$, bei 22
auf die $-$-Seite von $b$. Kehrt man zum Anfangspunkt (vor 1)
zurück, so schließt sich die $4^{\mathrm{te}}$ Figur, die $d$
heisst.
\par\Needspace{30\baselineskip}
Die Lage gegenüber den schon gezeichneten Kurven ergibt sich also
